\chapter{Introduction}\label{chap:intro}

\section*{At a glance}
\begin{itemize}
\item \lead{Goal} A machine-checked theorem about the shipping eraser behind \code{\#erase}, stated
  over lean4lean's formal model of Lean's kernel.
\item \lead{Status} \stOpen{} \hl{Nothing is verified}: the repository contains no verification
  declaration. The scope and the statements are being determined; checkpoint 1, the scope note,
  awaits the owner's approval.
\item \lead{Lean files} \code{Basic.lean}, \code{Erasure.lean} and \code{Printing.lean} under
  \code{LeanToLambdaBox/}: the shipping code.
\item \lead{Registers} \code{doc/SHIPPING-CHANGES.md} (Chapter~\ref{chap:shipping-changes}) and
  \code{doc/DIVERGENCES.md} (Chapter~\ref{chap:divergences}), rendered by a script.
\end{itemize}

\section{What is being verified}\label{sec:intro-what}

\begin{itemize}
\item \lead{The eraser} \code{lean-to-lambdabox} is the Lean~4 frontend of Peregrine. The command
  \code{\#erase t config c to "f.ast" mli "f.mli"} elaborates the term \code{t}, erases it and the
  constants it depends on to a $\lambda_\Box$ program, and writes three files:
  \begin{itemize}
  \item \code{f.ast}: the program as an S-expression, which Peregrine (\code{peregrine-tool})
        reads and compiles;
  \item \code{f.ast.inlinings}: the constants Peregrine should inline;
  \item \code{f.mli}: an OCaml signature for the compiled program.
  \end{itemize}
  Section~\ref{sec:trust-shipping} describes the code.
\item \lead{The goal} A theorem about this shipping code itself, built on lean4lean's model of
  Lean's kernel. It covers the inputs whose verification needs nothing beyond what lean4lean
  \code{master} provides; determining that scope is open (Chapter~\ref{chap:open}).
\item \lead{The references} The verification follows these as closely as the differences between
  Lean and Rocq allow; every divergence from the first two is an entry of
  \code{doc/DIVERGENCES.md}.
\end{itemize}

\begin{tabular}{p{0.3\linewidth}p{0.62\linewidth}}
\toprule
Reference & Relevant part \\
\midrule
Letouzey, \emph{A New Extraction for Coq}, TYPES 2002 & Section 3: the extraction function and
  its correctness (Theorems 12, 13, 15) \\
Sozeau et al., \emph{Correct and Complete Type Checking and Certified Erasure for Coq, in Coq},
  J. ACM 72(1), 2025 & Section 7: the target $\lambda_\Box$, the erasure function and relation,
  \code{erases\_correct}; the MetaRocq sources under \code{erasure/theories/} \\
Dima, \emph{Compiling Lean programs with Rocq's extraction pipeline}, MPRI report, 2025 & Section
  4: the design of this eraser \\
\bottomrule
\end{tabular}

\section{Pins}\label{sec:intro-pins}

The table is read from \code{lean-toolchain} and \code{lake-manifest.json} of the documented
commit. The \code{lean4lean} package is the fork \code{github.com/barabbs/lean4lean}, branch
\code{master}.

%% Generated by blueprint/scripts/render_registers.py from lean-toolchain, lake-manifest.json. Do not edit.
\begin{tabular}{lp{0.5\linewidth}l}
\toprule
Pin & Value & Source \\
\midrule
Lean toolchain & \code{leanprover/\allowbreak{}lean4:v4.33.0-rc2} & \code{lean-toolchain} \\
\code{lean4lean} & \code{https:/\allowbreak{}/\allowbreak{}github.com/\allowbreak{}barabbs/\allowbreak{}lean4lean} at \code{8223d223} (pinned by commit) & required by \code{lakefile.toml} \\
\code{batteries} & \code{https:/\allowbreak{}/\allowbreak{}github.com/\allowbreak{}leanprover-community/\allowbreak{}batteries} at \code{76e1c118} (input \code{v4.33.0-rc2}) & inherited from a dependency \\
\bottomrule
\end{tabular}


\section{Branches}\label{sec:intro-branches}

\begin{tabular}{llp{0.55\linewidth}}
\toprule
Branch & Merges from & Content \\
\midrule
\code{shipping} & -- & the shipping code on the lean4lean base, and the strictly necessary fixes;
  the register \code{doc/SHIPPING-CHANGES.md} \\
\code{verification} & \code{shipping} & the verification code and the register
  \code{doc/DIVERGENCES.md}; it never edits shipping files \\
\code{blueprint} & \code{verification} & this blueprint, under \code{blueprint/} \\
\bottomrule
\end{tabular}

\begin{itemize}
\item \lead{Merge flow} \code{shipping} to \code{verification} to \code{blueprint}, never
  backwards; every merge is \code{git merge -{}-no-ff}.
\item \lead{Documented commit} The blueprint describes the commit it is built from.
\end{itemize}

\section{Reading this blueprint}\label{sec:intro-reading}

\begin{tabular}{lp{0.7\linewidth}}
\toprule
Chapter & Content \\
\midrule
\ref{chap:trust} & The trust boundary: what is proved, the shipping code, what is inherited from
  lean4lean \\
\ref{chap:shipping-changes} & The shipping-change register, rendered \\
\ref{chap:divergences} & The divergence register, rendered \\
\ref{chap:open} & What is open \\
\bottomrule
\end{tabular}

\begin{itemize}
\item \lead{Nodes} Every definition and result cites its Lean declarations (\code{\textbackslash
  lean}). In the dependency graph, a green border means the declarations exist and depend on
  allowed axioms only.
\item \lead{Allowed axioms} \code{propext}, \code{Classical.choice}, \code{Quot.sound}, and the
  sorry sources and axioms of lean4lean; a node that depends on the latter lists them
  (\emph{Inherited from lean4lean}).
\item \lead{Badges} \stShipping{} shipping code, described, not verified; \stInherited{} trust
  inherited from lean4lean; \stOpen{} not done; \stProved{} proved.
\item \lead{Checks} \code{blueprint/scripts/audit.py} checks every node against the Lean
  environment and the generated chapters against their sources (\code{blueprint/README.md}).
\end{itemize}
